\documentclass[12pt,a4paper]{article}

% ============================================================
% Paquetes
% ============================================================
\usepackage[utf8]{inputenc}
\usepackage[T1]{fontenc}
\usepackage[spanish]{babel}
\usepackage[margin=2.5cm]{geometry}
\usepackage{graphicx}
\usepackage{hyperref}
\usepackage{listings}
\usepackage{xcolor}
\usepackage{booktabs}
\usepackage{enumitem}
\usepackage{fancyhdr}
\usepackage{titlesec}

% ============================================================
% Configuración
% ============================================================
\hypersetup{
    colorlinks=true,
    linkcolor=blue,
    urlcolor=blue,
    pdftitle={Manual de Usuario - KineticForge v1.0},
    pdfauthor={Equipo KineticForge}
}

\definecolor{codegray}{rgb}{0.5,0.5,0.5}
\definecolor{codegreen}{rgb}{0,0.6,0}
\definecolor{codepurple}{rgb}{0.58,0,0.82}

\lstdefinestyle{mystyle}{
    backgroundcolor=\color{gray!10},
    commentstyle=\color{codegreen},
    keywordstyle=\color{magenta},
    stringstyle=\color{codepurple},
    basicstyle=\ttfamily\small,
    breaklines=true,
    frame=single,
    numbers=left,
    numberstyle=\tiny\color{codegray}
}
\lstset{style=mystyle}

\pagestyle{fancy}
\fancyhf{}
\fancyhead[L]{KineticForge v1.0}
\fancyhead[R]{\thepage}
\renewcommand{\headrulewidth}{0.4pt}

% ============================================================
% Documento
% ============================================================
\begin{document}

\begin{titlepage}
    \centering
    \vspace*{3cm}

    {\Huge\bfseries KineticForge\par}
    \vspace{0.5cm}
    {\Large Manual de Usuario\par}
    \vspace{0.3cm}
    {\large Versión 1.0.0\par}

    \vspace{3cm}

    {\large Convierte hojas A4 con dibujos a mano\\en GIF animados con fondo transparente\par}

    \vfill

    {\large \today\par}
\end{titlepage}

\tableofcontents
\newpage

% ============================================================
\section{Introducción}
% ============================================================

\textbf{KineticForge} es una aplicación de escritorio que automatiza la conversión de
hojas A4 con dibujos a mano (organizados en una grilla) en \textbf{GIF animados con
fondo transparente}.

\subsection{¿Para quién es?}

\begin{itemize}[leftmargin=*]
    \item Dibujantes tradicionales que quieren digitalizar sus animaciones
    \item Estudiantes de animación que necesitan GIFs para portafolio
    \item Docentes que enseñan animación tradicional
    \item Cualquier persona con un escáner y ganas de animar
\end{itemize}

\subsection{¿Qué resuelve?}

Los flujos tradicionales requieren:
\begin{itemize}[leftmargin=*]
    \item Photoshop (suscripción mensual)
    \item Recortar 48 frames a mano (30-60 min)
    \item Alinear los frames para evitar vibración
    \item Exportar con transparencia
\end{itemize}

\textbf{KineticForge hace todo eso en menos de 2 minutos.}

% ============================================================
\section{Requisitos}
% ============================================================

\subsection{Hardware}
\begin{itemize}[leftmargin=*]
    \item \textbf{Sistema operativo:} Windows 10+, Linux (Ubuntu 20.04+), macOS 11+
    \item \textbf{RAM:} 4 GB mínimo, 8 GB recomendado
    \item \textbf{Disco:} 500 MB libres
    \item \textbf{Escáner} (o cámara de celular) para digitalizar la plantilla
\end{itemize}

\subsection{Software}
\begin{itemize}[leftmargin=*]
    \item \textbf{Java 21} (para ejecutar desde JAR)
    \item \textbf{O instalador nativo} (que incluye Java — recomendado)
\end{itemize}

% ============================================================
\section{Instalación}
% ============================================================

\subsection{Opción A: Instalador nativo (recomendado)}

\textbf{Linux (.deb):}
\begin{lstlisting}[language=bash]
sudo dpkg -i kineticforge-1.0.0.deb
\end{lstlisting}

\textbf{Windows (.exe):}
\begin{itemize}[leftmargin=*]
    \item Doble clic en \texttt{kineticforge-1.0.0.exe}
    \item Seguir el asistente
\end{itemize}

\textbf{macOS (.dmg):}
\begin{itemize}[leftmargin=*]
    \item Doble clic en \texttt{kineticforge-1.0.0.dmg}
    \item Arrastrar a la carpeta Aplicaciones
\end{itemize}

\subsection{Opción B: Ejecutar desde JAR}

Requiere Java 21 instalado.

\begin{lstlisting}[language=bash]
java -jar kineticforge-1.0.0.jar
\end{lstlisting}

% ============================================================
\section{Primeros pasos}
% ============================================================

\subsection{Generar la plantilla}

KineticForge incluye un generador de plantillas. \textbf{Antes de dibujar, necesitás
la plantilla correcta.}

\textbf{Desde la línea de comandos:}
\begin{lstlisting}[language=bash]
java -jar kineticforge-1.0.0.jar --generate-template --preset compacta
\end{lstlisting}

\textbf{O desde la app:}
\begin{enumerate}[leftmargin=*]
    \item Abrir KineticForge
    \item Ir a la sección ``Plantillas''
    \item Elegir preset y click en ``Generar''
    \item Se guarda en \texttt{\textasciitilde/kineticforge/plantillas/}
\end{enumerate}

\subsection{Imprimir la plantilla}
\begin{itemize}[leftmargin=*]
    \item \textbf{Papel:} A4 (210 × 297 mm)
    \item \textbf{Calidad:} alta (300 DPI mínimo)
    \item \textbf{Márgenes:} los que tiene la plantilla (no ajustar)
    \item \textbf{No usar ``ajustar a página''} — imprimir a tamaño real
\end{itemize}

\subsection{Dibujar}
\begin{itemize}[leftmargin=*]
    \item Dibujar \textbf{un cuadro por celda} (el orden es de izquierda a derecha, arriba a abajo)
    \item \textbf{El contorno del personaje debe estar cerrado} (importante para la detección automática)
    \item Dibujar en \textbf{tinta negra} o \textbf{lápiz oscuro}
\end{itemize}

\subsection{Escanear}
\begin{itemize}[leftmargin=*]
    \item \textbf{Resolución:} 300 DPI mínimo
    \item \textbf{Formato:} PNG (sin pérdida) o JPG (calidad alta)
    \item \textbf{No recortar} (dejar la hoja completa)
\end{itemize}

% ============================================================
\section{Guía del wizard}
% ============================================================

KineticForge usa un wizard de 4 pasos.

\subsection{Paso 1: Cargar imagen}
\begin{itemize}[leftmargin=*]
    \item Arrastrar la imagen escaneada al área de carga
    \item O click en ``Elegir archivo...''
    \item Formatos aceptados: PNG, JPG, TIFF, BMP
\end{itemize}

\textbf{Validación:} la app muestra el nombre del archivo y sus dimensiones.

\subsection{Paso 2: Detectar grilla}
\begin{itemize}[leftmargin=*]
    \item Elegir preset (Compacta, Densa, etc.)
    \item Elegir orientación (Vertical / Horizontal)
    \item Ajustar margen si es necesario
    \item Click en ``Detectar automáticamente''
\end{itemize}

\textbf{Resultado:} la app muestra las celdas detectadas con líneas azules.

\textbf{Si la detección falla:} ajustar preset, orientación o margen manualmente.

\subsection{Paso 3: Procesar frames}
\begin{itemize}[leftmargin=*]
    \item Click en ``Procesar frames''
    \item La app:
    \begin{enumerate}[leftmargin=*]
        \item Extrae cada celda como un frame
        \item Quita el fondo blanco
        \item Alinea los frames para evitar vibración
    \end{enumerate}
\end{itemize}

\textbf{Duración:} \textasciitilde 10-30 segundos según la resolución.

\textbf{Resultado:} preview con las 48 miniaturas.

\subsection{Paso 4: Exportar}
\begin{itemize}[leftmargin=*]
    \item Elegir formato (GIF, PNG spritesheet, WebP)
    \item Ajustar FPS (1-30)
    \item Elegir si reproducir en loop
    \item Click en ``Exportar...''
\end{itemize}

\textbf{Resultado:} el archivo se guarda en \texttt{\textasciitilde/kineticforge/exportaciones/}.

% ============================================================
\section{Formatos de la plantilla}
% ============================================================

KineticForge ofrece 6 presets:

\begin{table}[h]
\centering
\begin{tabular}{@{}lcccc@{}}
\toprule
\textbf{Preset} & \textbf{Grid} & \textbf{Total} & \textbf{Uso} \\
\midrule
Detalle       & 4×3   & 12  & Bocetos grandes \\
Estándar      & 6×4   & 24  & 1 s a 24 fps \\
Compacta      & 8×6   & 48  & 2 s a 24 fps \\
Densa         & 12×8  & 96  & 4 s a 24 fps \\
Cinemática    & 16×6  & 96  & Widescreen \\
Súper densa   & 24×12 & 288 & 12 s a 24 fps \\
\bottomrule
\end{tabular}
\end{table}

% ============================================================
\section{Exportación}
% ============================================================

\subsection{Formatos disponibles}

\begin{table}[h]
\centering
\begin{tabular}{@{}lccc@{}}
\toprule
\textbf{Formato} & \textbf{Transparencia} & \textbf{Colores} & \textbf{Uso} \\
\midrule
GIF              & Sí (1 bit)   & 256   & Universal \\
PNG spritesheet  & Sí (8 bits)  & 16M   & Motores de juego \\
WebP             & Sí (8 bits)  & 16M   & Web moderna \\
\bottomrule
\end{tabular}
\end{table}

\subsection{Parámetros}
\begin{itemize}[leftmargin=*]
    \item \textbf{FPS:} 1 a 30 fotogramas por segundo
    \item \textbf{Loop:} reproducir infinitamente o una sola vez
\end{itemize}

% ============================================================
\section{Preguntas frecuentes (FAQ)}
% ============================================================

\subsection{¿Puedo usar una foto de celular en vez de escáner?}

Sí. Asegurate de:
\begin{itemize}[leftmargin=*]
    \item Buena iluminación (sin sombras)
    \item Hoja plana (sin dobleces)
    \item Resolución alta ($>$2000 px de ancho)
\end{itemize}

\subsection{¿Qué pasa si el personaje toca el borde de la celda?}

El \texttt{BackgroundRemover} puede recortar el contorno. \textbf{Dejá al menos 2 mm de margen}
entre el personaje y el borde de la celda.

\subsection{¿Por qué el GIF se ve con ``banding''?}

Es limitación del formato GIF (256 colores). Para mejor calidad, exportá a WebP
o PNG spritesheet.

\subsection{¿Puedo animar más de 48 frames?}

Sí, elegí un preset con más celdas (Densa, Súper densa). También podés combinar
varias hojas (feature pendiente para v1.1).

% ============================================================
\section{Solución de problemas}
% ============================================================

\subsection{La detección de grilla falla}

\textbf{Síntomas:} las celdas no se alinean con la plantilla.

\textbf{Soluciones:}
\begin{enumerate}[leftmargin=*]
    \item Verificar que la imagen sea de la plantilla (no una hoja en blanco)
    \item Verificar orientación (Vertical vs Horizontal)
    \item Ajustar el preset (Compacta, Densa, etc.)
    \item Ajustar el margen (slider)
\end{enumerate}

\subsection{Los frames salen vacíos}

\textbf{Síntomas:} las miniaturas están todas blancas.

\textbf{Causa:} el \texttt{BackgroundRemover} borra demasiado.

\textbf{Solución:} verificar que la plantilla no tenga fondo gris.

\subsection{El GIF pesa mucho}

\textbf{Solución:} reducir FPS, usar menos colores, o exportar a WebP (más eficiente).

\subsection{Error ``Java no encontrado''}

\textbf{Solución:} instalar Java 21 desde \url{https://adoptium.net}.

% ============================================================
\section{Detalles técnicos}
% ============================================================

\subsection{Arquitectura}

KineticForge usa \textbf{arquitectura por capas:}

\begin{verbatim}
UI (JavaFX)
   |
Controller / ViewModel
   |
Core (logica pura)
   |
IO (carga, exporta, persiste)
\end{verbatim}

\subsection{Pipeline de procesamiento}

\begin{enumerate}[leftmargin=*]
    \item \textbf{Detección de grilla:} proyección de histogramas + filtro de líneas espaciadas uniformemente
    \item \textbf{Extracción de frames:} recorte de cada celda
    \item \textbf{Eliminación de fondo:} flood fill desde los 4 bordes con antialiasing
    \item \textbf{Alineación:} bounding box en canvas común
    \item \textbf{Exportación:} GIF (javax.imageio), PNG spritesheet, WebP
\end{enumerate}

\subsection{Tecnologías}
\begin{itemize}[leftmargin=*]
    \item \textbf{Java 21} (LTS)
    \item \textbf{JavaFX 21} (UI)
    \item \textbf{Apache PDFBox 3} (plantillas PDF)
    \item \textbf{TwelveMonkeys ImageIO} (TIFF, BMP)
    \item \textbf{AtlantaFX} (temas visuales)
    \item \textbf{JUnit 5} (testing — 157 tests)
    \item \textbf{Maven} (build)
\end{itemize}

\subsection{Licencia}

Apache License 2.0.

% ============================================================
\section{Contacto}
% ============================================================

\begin{itemize}[leftmargin=*]
    \item \textbf{Repo:} \url{https://github.com/2023207010est-blip/kineticforge}
    \item \textbf{Issues:} \url{https://github.com/2023207010est-blip/kineticforge/issues}
\end{itemize}

\end{document}
